% Part: computability
% Chapter: recursive-functions
% Section: pr-functions-computable

\documentclass[../../../include/open-logic-section]{subfiles}

\begin{document}

\olfileid{cmp}{rec}{cmp}
\olsection{الدوال العودية البدائية قابلة للحساب}

لتكن الدالة $\OLId{latin-ordinary}{h}$ معرّفة بالعودية البدائية كما يأتي:
\begin{eqnarray*}
\OLId{latin-ordinary}{h}(\vec \OLId{latin-ordinary}{x}, \OLMathNumeral{0})     & = & \OLId{latin-ordinary}{f}(\vec \OLId{latin-ordinary}{x}) \\
\OLId{latin-ordinary}{h}(\vec \OLId{latin-ordinary}{x}, \OLId{latin-ordinary}{y} + \OLMathNumeral{1}) & = & \OLId{latin-ordinary}{g}(\vec \OLId{latin-ordinary}{x}, \OLId{latin-ordinary}{y}, \OLId{latin-ordinary}{h}(\vec \OLId{latin-ordinary}{x}, \OLId{latin-ordinary}{y}))
\end{eqnarray*}
ولنفرض قابلية $\OLId{latin-ordinary}{f}$ و$\OLId{latin-ordinary}{g}$ للحساب. والرمز $\vec \OLId{latin-ordinary}{x}$ عندنا اختصار
لـ$\OLId{latin-ordinary}{x}_\OLMathNumeral{0}$، …،~$\OLId{latin-ordinary}{x}_{\OLId{latin-ordinary}{k}-\OLMathNumeral{1}}$. أما $\OLId{latin-ordinary}{h}(\vec \OLId{latin-ordinary}{x},\OLMathNumeral{0})$ فيتأتى حسابها
لأنها عين $\OLId{latin-ordinary}{f}(\vec \OLId{latin-ordinary}{x})$، وهي قابلة للحساب بالفرض. ثم يتأتى
حساب $\OLId{latin-ordinary}{h}(\vec \OLId{latin-ordinary}{x},\OLMathNumeral{1})$؛ إذ $\OLMathNumeral{1}=\OLMathNumeral{0}+\OLMathNumeral{1}$، فتكون $\OLId{latin-ordinary}{h}(\vec \OLId{latin-ordinary}{x},\OLMathNumeral{1})$
\begin{align*}
 \OLId{latin-ordinary}{h}(\vec \OLId{latin-ordinary}{x}, \OLMathNumeral{1}) & = \OLId{latin-ordinary}{g}(\vec \OLId{latin-ordinary}{x}, \OLMathNumeral{0}, \OLId{latin-ordinary}{h}(\vec \OLId{latin-ordinary}{x}, \OLMathNumeral{0})) =  \OLId{latin-ordinary}{g}(\vec \OLId{latin-ordinary}{x}, \OLMathNumeral{0}, \OLId{latin-ordinary}{f}(\vec \OLId{latin-ordinary}{x})).
\intertext{ويمكننا المواصلة على هذا النحو وحساب}
\OLId{latin-ordinary}{h}(\vec \OLId{latin-ordinary}{x}, \OLMathNumeral{2}) & = \OLId{latin-ordinary}{g}(\vec \OLId{latin-ordinary}{x}, \OLMathNumeral{1}, \OLId{latin-ordinary}{h}(\vec \OLId{latin-ordinary}{x}, \OLMathNumeral{1})) = \OLId{latin-ordinary}{g}(\vec \OLId{latin-ordinary}{x}, \OLMathNumeral{1}, \OLId{latin-ordinary}{g}(\vec \OLId{latin-ordinary}{x}, \OLMathNumeral{0}, \OLId{latin-ordinary}{f}(\vec \OLId{latin-ordinary}{x})))\\
\OLId{latin-ordinary}{h}(\vec \OLId{latin-ordinary}{x}, \OLMathNumeral{3}) & = \OLId{latin-ordinary}{g}(\vec \OLId{latin-ordinary}{x}, \OLMathNumeral{2}, \OLId{latin-ordinary}{h}(\vec \OLId{latin-ordinary}{x}, \OLMathNumeral{2})) = \OLId{latin-ordinary}{g}(\vec \OLId{latin-ordinary}{x}, \OLMathNumeral{2}, \OLId{latin-ordinary}{g}(\vec \OLId{latin-ordinary}{x}, \OLMathNumeral{1}, \OLId{latin-ordinary}{g}(\vec \OLId{latin-ordinary}{x}, \OLMathNumeral{0}, \OLId{latin-ordinary}{f}(\vec \OLId{latin-ordinary}{x}))))\\
\OLId{latin-ordinary}{h}(\vec \OLId{latin-ordinary}{x}, \OLMathNumeral{4}) & = \OLId{latin-ordinary}{g}(\vec \OLId{latin-ordinary}{x}, \OLMathNumeral{3}, \OLId{latin-ordinary}{h}(\vec \OLId{latin-ordinary}{x}, \OLMathNumeral{3})) = \OLId{latin-ordinary}{g}(\vec \OLId{latin-ordinary}{x}, \OLMathNumeral{3}, \OLId{latin-ordinary}{g}(\vec \OLId{latin-ordinary}{x}, \OLMathNumeral{2}, \OLId{latin-ordinary}{g}(\vec \OLId{latin-ordinary}{x}, \OLMathNumeral{1}, \OLId{latin-ordinary}{g}(\vec \OLId{latin-ordinary}{x}, \OLMathNumeral{0}, \OLId{latin-ordinary}{f}(\vec \OLId{latin-ordinary}{x})))))\\
& \vdots
\end{align*}
فإذا طلبنا $\OLId{latin-ordinary}{h}(\vec \OLId{latin-ordinary}{x},\OLId{latin-ordinary}{y})$ لأي حجة، حسبنا على التعاقب $\OLId{latin-ordinary}{h}(\vec \OLId{latin-ordinary}{x},\OLMathNumeral{0})$ و
$\OLId{latin-ordinary}{h}(\vec \OLId{latin-ordinary}{x},\OLMathNumeral{1})$ و…، حتى نبلغ $\OLId{latin-ordinary}{h}(\vec \OLId{latin-ordinary}{x},\OLId{latin-ordinary}{y})$.

فقد بان أن التعريف بالعودية البدائية يحفظ قابلية الحساب متى كانت
الدالتان $\OLId{latin-ordinary}{f}$ و$\OLId{latin-ordinary}{g}$ قابلتين للحساب. وكذلك التركيب يحفظ هذه
القابلية إذا تحققت في $\OLId{latin-ordinary}{f}$ وفي الدوال $\OLId{latin-ordinary}{g}_\OLId{latin-ordinary}{i}$.

وحيث إن الدوال الأساسية $\Zero$ و$\Succ$ و$\Proj{\OLId{latin-ordinary}{n}}{\OLId{latin-ordinary}{i}}$ قابلة
للحساب، وإن التركيب والعودية البدائية لا يخرجان بنا من
الدوال القابلة للحساب، لزم أن تكون كل دالة عودية بدائية قابلة للحساب.

\end{document}
